% A homogeneous space M: the group acts transitively, so the one point a is carried to two
% different targets b and b' by two different group elements g and g'.
\documentclass{standalone}
\usepackage{geometricfigures}
\usepackage{pgfplots}
\pgfplotsset{compat=1.18}

\begin{document}
\begin{tikzpicture}
    \begin{axis}[
        xlabel={$x$},
        ylabel={$y$},
        zlabel={$z$},
        grid=major,
        view={-30}{30}, % Adjust the view for better visibility
        colormap/viridis, % Use the viridis colormap
    ]
        % (ii) Plot the 2D manifold with the Julia-like (viridis) colormap
        \addplot3[
            surf,
            shader=interp,
            domain=-5:5,
            y domain=-5:5,
            samples=50,
            opacity=0.9, % Make the surface slightly transparent
        ]{sin(deg(.5*x)) * cos(deg(.5*y))};

        % ---- the orbit of a: two targets, two group elements ------------------
        % g . a = b,   along (-4,4) -> (4,-4)
        \addplot3[
            domain=0:1, samples=100, samples y=0, color=mred, thick, opacity=.9,
        ]({-4 + 8*x}, {4 - 8*x}, {sin(deg(.5*(-4 + 8*x))) * cos(deg(.5*(4 - 8*x)))});
        % g' . a = b',  along (-4,4) -> (0,-4).  Both targets are on the front edge
        % y = -4, which is the only edge this view (-30/30) does not fold into the
        % middle of the picture: the corners (-4,-4) and (4,4) project inside the
        % surface and a point put there lands on top of the other curve.
        \addplot3[
            domain=0:1, samples=100, samples y=0, color=mpurple, thick, opacity=.9,
        ]({-4 + 4*x}, {4 - 8*x}, {sin(deg(.5*(-4 + 4*x))) * cos(deg(.5*(4 - 8*x)))});

        % (i) the source point and the two images
        \addplot3[only marks, mark=*, mark size=3pt, color=mred] coordinates {
            (-4, 4, {sin(deg(-2)) * cos(deg(2))})   % a
            (4, -4, {sin(deg(2)) * cos(deg(-2))})   % b
        };
        \addplot3[only marks, mark=*, mark size=3pt, color=mpurple] coordinates {
            (0, -4, {0 * cos(deg(-2))})             % b'
        };

        \node[above right, color=mred]    at (axis cs:-4,4,{sin(deg(-2)) * cos(deg(2))}) {$a$};
        \node[above right, color=mred]    at (axis cs:4,-4,{sin(deg(2)) * cos(deg(-2))}) {$b$};
        \node[below right, color=mpurple] at (axis cs:0,-4,0) {$b'$};
        \node[above right, color=fg]   at (axis cs:4,-4,-.8) {$\mathcal{M}$};

        % each curve carries the element that realises it
        \node[color=mred]    at (axis cs: 1.0,-1.0,.95)   {${\color{mred}g}$};
        \node[color=mpurple] at (axis cs:-3.3,-1.0,-.95) {${\color{mpurple}g'}$};
    \end{axis}
\end{tikzpicture}
\end{document}
